\documentclass[12pt,a4paper]{article}
\usepackage[margin=1in]{geometry}
\usepackage{graphicx}
\usepackage{float}
\usepackage{booktabs}
\usepackage{hyperref}
\usepackage{longtable}

\title{CA5 Question 2 Report\\TravBot Live Agent System}
\author{tahamajs\\Student ID: 810101504}
\date{\today}

\newcommand{\safeimage}[3]{
\begin{figure}[H]
\centering
\IfFileExists{#1}{\includegraphics[width=#2\textwidth]{#1}}{\fbox{\texttt{Missing figure: \detokenize{#1}}}}
\caption{#3}
\end{figure}
}

\begin{document}
\maketitle

\section{Summary}
This report documents the final Q2 implementation of TravBot as a strict live-integration multi-tool travel agent. The system integrates OpenAI, Amadeus, Tavily, Kaggle datasets, LlamaParse, LanceDB, and LangGraph. The code is modularized under \texttt{answers/Q2/src} and produces reproducible scenario logs and report plots.

\section{Implementation Outputs}
\begin{longtable}{p{0.45\textwidth} p{0.45\textwidth}}
\toprule
Artifact & Description \\
\midrule
\texttt{answers/Q2/lancedb\_data/} & Vector database for FAQ and tourism content \\
\texttt{answers/Q2/results/scenario\_results.csv} & Scenario-level outputs, latency, and tool usage \\
\texttt{answers/Q2/results/scenario\_results.json} & JSON export of scenario outputs \\
\texttt{answers/reports/figures/q2/*.png} & Report-ready scenario and performance plots \\
\bottomrule
\end{longtable}

\section{System Architecture}
The Q2 system follows this deterministic sequence:
\begin{enumerate}
\item Runtime key loading from environment files
\item Live preflight checks for tools, modules, and credentials
\item Kaggle dataset ingestion for IATA and currency mappings
\item LlamaParse extraction from the tourism PDF
\item LanceDB initialization for FAQ and tourism retrieval
\item LangGraph ReAct agent construction with seven tools
\item Scenario execution and output logging
\item Plot generation for report publication
\end{enumerate}

\section{Tooling Coverage}
The agent uses the following tools:
\begin{itemize}
\item Flight search via Amadeus
\item Hotel search via Amadeus
\item Restaurant search via Tavily
\item Weather information via Tavily
\item Currency conversion context via Tavily and Kaggle mappings
\item FAQ semantic search via LanceDB
\item Trip planning synthesis using OpenAI and retrieved contexts
\end{itemize}

\section{Figures}
\safeimage{figures/q2/q2_scenario_latency.png}{0.88}{Scenario latency by scenario ID}
\safeimage{figures/q2/q2_latency_distribution.png}{0.78}{Latency distribution}
\safeimage{figures/q2/q2_tool_usage_frequency.png}{0.88}{Tool usage frequency}
\safeimage{figures/q2/q2_query_length_vs_latency.png}{0.72}{Query length vs latency}

\section{Reproducible Execution}
All commands are executed under:
\begin{verbatim}
source /Users/tahamajs/Documents/uni/venv/bin/activate
\end{verbatim}

Primary command:
\begin{verbatim}
bash answers/scripts/run_all.sh
\end{verbatim}

\section{Validation Criteria}
The final implementation targets the following checks:
\begin{itemize}
\item Strict live preflight passes with required credentials and services.
\item Notebook execution completes from top to bottom.
\item Scenario output files are generated and non-empty.
\item All Q2 figures are generated under \texttt{answers/reports/figures/q2}.
\item Report builds succeed with XeLaTeX into \texttt{answers/reports/build}.
\end{itemize}

\end{document}
